\chapter{State of the Art}
\label{sec:stateoftheart}

The detection of fraud in financial transactions is a challenging problem at the intersection of machine learning, class imbalance, and operational deployment constraints. This chapter surveys five areas that define the present study: the fraud detection problem and its benchmark datasets, evaluation under extreme class imbalance, classical machine learning, deep learning and transformer-based models for tabular data, and explainability. It concludes by synthesising recurring methodological limitations and identifying the research gaps that motivate the experimental design.

% ============================================================
\section{Fraud Detection in Financial Transactions}
\label{sec:sota-fraud-detection-financial}
% ============================================================

\subsection{The Financial Fraud Problem}
\label{subsec:sota-fraud-nature}

Financial fraud in electronic payment systems causes significant economic harm, with annual global losses consistently estimated in the hundreds of billions of dollars \cite{rymantubb2018survey}. The term encompasses a broad range of adversarial behaviours — card-not-present fraud, account takeover, synthetic identity fraud, and money laundering — each exhibiting distinct statistical signatures in transaction data \cite{hilal2022financial}. What unites these phenomena from a machine learning perspective is a combination of extreme class imbalance, concept drift as fraudsters adapt to detection systems, and the asymmetric cost of errors: a missed fraud (false negative) is typically far more costly than a legitimate transaction incorrectly flagged (false positive) \cite{ali2022financial}.

Rule-based expert systems were historically the dominant detection paradigm, offering interpretability and direct encoding of domain knowledge at the cost of continuing manual maintenance and limited generalisation to novel attack patterns \cite{rymantubb2018survey}. Supervised machine learning uses historical labelled transactions to model feature interactions and can be updated as new labelled data become available without requiring every decision rule to be specified explicitly. Its practical limitation is data availability: fraud labels are costly to produce, may arrive with delay, and the underlying distribution shifts as fraudsters respond to deployed models \cite{hernandezaros2024financial}. These constraints motivate the use of benchmark datasets to standardise evaluation and the assessment of model behaviour under distribution shift.

\subsection{Benchmark Datasets}
\label{subsec:sota-datasets}

Reproducible comparative studies benefit from publicly available datasets with documented characteristics. The following two benchmark collections are central to this thesis.

The \textbf{ULB Credit Card Fraud dataset} \cite{dalpozzolo2017realistic} comprises 284\,807 transactions made by European cardholders over two days in September 2013, with 492 fraudulent transactions (fraud rate 0.172\%). Sensitive features were transformed through principal component analysis (PCA), yielding 28 anonymised numerical components (V1--V28) alongside the elapsed \texttt{Time} and transaction \texttt{Amount}. The published representation limits disclosure of the original variables while retaining a numerical feature space for modelling. Its principal limitations are temporal and semantic: the two-day window constrains temporal generalisation studies, and the transformed features preclude domain-specific interpretation and feature engineering. Despite these constraints, ULB remains a widely used benchmark in academic fraud detection research \cite{ali2022financial, mienye2024deep}.

The \textbf{Bank Account Fraud (BAF)} dataset suite \cite{jesus2022turning}, introduced at NeurIPS 2022, provides a synthetic yet realistic alternative designed to overcome several limitations of ULB. BAF Base contains one million bank account opening applications with thirty modelling features — twenty-five numerical and five categorical (\texttt{payment\_type}, \texttt{employment\_status}, \texttt{housing\_status}, \texttt{source}, \texttt{device\_os}) — and a fraud rate of approximately 1.1\%. Unlike ULB, BAF preserves semantically meaningful feature names, enabling interpretation of model decisions in domain terms. The suite includes five controlled variants: Variant~I increases group-size disparity while equalising group prevalence; Variant~II increases group-wise prevalence disparity while equalising group sizes; Variant~III introduces group-wise separability differences through two synthetic features; Variant~IV changes prevalence disparity over time; and Variant~V changes separability over time through the same synthetic features. This structure supports evaluation under controlled bias and distributional conditions without requiring access to production environments.

These two datasets span different problem regimes. ULB is a dense, low-dimensional binary classification problem with extreme imbalance in a temporally constrained snapshot. BAF is a large-scale, semantically rich multi-variant benchmark designed to expose robustness limitations. Evaluating the same model families with a shared outer holdout and reporting protocol across both datasets reduces dependence on any single data regime, as argued in the context of multi-dataset fraud benchmarks \cite{grover2022fdb}.

% ============================================================
\section{Evaluation of Fraud Detection Systems under Class Imbalance}
\label{subsec:fraud_detection_class_imbalance}
% ============================================================

\subsection{The Accuracy Paradox and Metric Selection}
\label{subsec:sota-accuracy-paradox}

Under extreme class imbalance — where fraud constitutes 0.1\%--2\% of all transactions — a classifier that predicts every transaction as legitimate achieves accuracy above 98\% without detecting a single fraud. This \textit{accuracy paradox} renders overall accuracy uninformative and motivates the use of metrics that directly measure performance on the minority class \cite{hegarcia2009learning}.

The area under the receiver operating characteristic curve (ROC-AUC) is widely reported in the fraud detection literature. It measures the probability that a randomly selected positive instance is ranked above a randomly selected negative instance and is invariant to class prevalence for fixed class-conditional score distributions. Under extreme imbalance, however, a small false-positive rate can still correspond to a large absolute number of false alerts, which the ROC curve does not expose directly \cite{saito2015precision}. Evidence from ULB illustrates the resulting interpretive risk: classifiers can exceed 0.97 ROC-AUC while exhibiting substantially weaker precision--recall performance \cite{imani2026why}. ROC-AUC is therefore informative as a ranking measure but can be insufficient as the primary operational summary for fraud detection.

The precision--recall curve directly represents the trade-off between precision and recall for the positive class, making its summary area an appropriate primary measure for imbalanced classification \cite{saito2015precision}. Average precision, the implementation used in this thesis and reported as PR-AUC, has a random-ranking baseline equal to class prevalence (e.g., approximately 0.011 on BAF Base). This prevalence-dependent reference is absent from ROC-AUC. The thesis therefore uses PR-AUC as its primary metric, complemented by $F_1$ and the recall-weighted $F_2$ score at a validation-selected threshold.

\subsection{Threshold Selection and Probability Calibration}
\label{subsec:sota-threshold}

Point metrics — precision, recall, $F_1$, and $F_{\beta}$ — are defined for a specific decision threshold $\tau$ applied to the model's output scores. For a calibrated posterior probability, $\tau=0.5$ corresponds to equal false-positive and false-negative costs; it is not generally appropriate when operational costs are asymmetric or probabilities are imperfectly calibrated. Selecting $\tau$ on the test set constitutes evaluation bias; it must be chosen on a held-out validation set strictly separated from the test data \cite{luzio2024decoupling}.

Multiple threshold selection strategies have been proposed. Fixing $\tau = 0.5$ is simple but frequently suboptimal. Data-driven alternatives include selecting the threshold that maximises $F_1$ (equal weight on precision and recall), maximising $F_2$ (with recall receiving four times the weight of precision through $\beta^2$), or constraining precision to a minimum operationally acceptable level before maximising recall \cite{dalpozzolo2015calibrating}. A validation score threshold can be selected without calibrated probabilities, but probability calibration is required when the score is interpreted as posterior risk or when a probability threshold is derived from explicit costs. Under-sampling strategies that alter the class ratio during training can distort posterior probability estimates and may require post-hoc calibration correction \cite{dalpozzolo2015calibrating}. The interaction between threshold strategy, calibration, and model architecture is an understudied aspect of the fraud detection literature.

\subsection{Class Imbalance Correction Strategies}
\label{subsec:sota-imbalance-strategies}

A large body of work addresses class imbalance at the data level. \textit{Synthetic Minority Oversampling Technique} (SMOTE) \cite{chawla2002smote} is the most widely adopted oversampling method, generating synthetic minority-class instances as convex combinations of existing ones in feature space. Extensions such as SMOTE+Tomek and SMOTEENN \cite{batista2004balancing} add boundary cleaning steps to remove ambiguous majority instances after oversampling, seeking a cleaner decision boundary.

Random undersampling (RUS) discards majority-class instances to achieve a target ratio. This is computationally efficient but discards potentially informative negative examples. Class-weight adjustment, wherein the loss function assigns greater importance to minority-class errors, avoids discarding data and is a common alternative when the majority class contains non-trivial patterns.

Survey evidence \cite{hegarcia2009learning, krawczyk2016learning} cautions that no single strategy consistently dominates across all datasets and model families. A critical implementation constraint is that oversampling must be applied exclusively within the training fold during cross-validation: applying it before the split introduces data leakage that can inflate reported performance \cite{hayat2025leakage}. The interaction between imbalance strategy, model family, and input representation is therefore an empirical question addressed by the factorial design of this thesis.

% ============================================================
\section{Classical Machine Learning for Fraud Detection}
\label{subsec:classicML}
% ============================================================

\subsection{Linear Models and Ensemble Methods}
\label{subsec:sota-linear-ensemble}

Logistic regression remains a standard baseline in fraud detection research. Its linear decision boundary, direct interpretability of coefficients, and probabilistic output make it a natural starting point for practitioners \cite{dalpozzolo2017realistic}; as with other probabilistic classifiers, its calibration should be checked empirically. Its limitation is expressive capacity: without engineered interaction or nonlinear terms, it cannot represent the nonlinear interaction patterns that may characterise sophisticated fraud.

Random forests construct an ensemble of decorrelated decision trees via bootstrap aggregation and random feature subsets. Averaging reduces variance relative to an individual tree, but does not by itself guarantee robustness to noise or distribution shift. Random forests can also be biased towards the majority class under severe imbalance, and their independently fitted trees do not use the sequential error correction employed by gradient boosting.

\subsection{Gradient-Boosted Decision Trees}
\label{subsec:sota-gbdt}

Gradient-boosted decision trees (GBDTs) have repeatedly achieved strong performance on tabular data \cite{shwartzziv2022tabular}. The gradient boosting framework constructs an additive ensemble of trees, each fitted to the pseudo-residuals of the current ensemble, so that successive weak learners address errors remaining from earlier stages.

LightGBM \cite{ke2017lightgbm} introduces leaf-wise tree growth and histogram-based approximate gradient computation to improve computational efficiency at scale. CatBoost \cite{prokhorenkova2018catboost} adds ordered boosting — a sequential scheme designed to reduce prediction-shift bias during training — alongside a specialised categorical feature encoder designed to limit target leakage within tree construction. These properties are relevant to financial datasets containing categorical fields such as device type or employment status.

Comparative studies on the ULB dataset frequently rank CatBoost and LightGBM among the strongest classical models \cite{leevy2023binary, singh2025benchmarking, thimonier2024comparative}. Their practical advantages include limited sensitivity to feature scale, implementation-level support for missing values, and computationally efficient exact feature attribution via TreeSHAP \cite{lundberg2020local}. Their probability calibration remains data- and configuration-dependent and should be evaluated rather than assumed.

\subsection{Cost-Sensitive and Anomaly Detection Approaches}
\label{subsec:sota-anomaly}

Cost-sensitive learning \cite{hoppner2021idcs} offers an alternative to resampling by directly incorporating the asymmetric misclassification costs into the model's objective function. Rather than modifying the training data distribution, cost-sensitive methods specify a loss matrix and optimise expected cost rather than expected error rate. This approach avoids the distributional distortion introduced by oversampling and is particularly natural for fraud detection, where the cost ratio between false negatives and false positives can be quantified from operational data.

One-class classification treats fraud detection as an anomaly detection problem: a model trained exclusively on legitimate transactions flags deviations as potential fraud. The one-class support vector machine (OCSVM) \cite{hegarcia2009learning} maps observations into a kernel feature space and learns a maximum-margin boundary that separates most training observations from the origin. These approaches are appealing when labelled fraud examples are scarce, but, when reliable labels are available, supervised methods can exploit information that a one-class objective deliberately excludes \cite{leevy2023binary, thimonier2024comparative}.

% ============================================================
\section{Deep Learning for Fraud Detection}
\label{subsec:deep_learning}
% ============================================================

Deep learning methods have been applied to fraud detection across a range of architectures, with uneven results. A recent review \cite{mienye2024deep} surveys the deployment of convolutional networks, recurrent models, and autoencoders across credit card fraud benchmarks.

Fully connected networks (MLPs) offer greater nonlinear capacity than logistic regression but require regularisation and hyperparameter tuning to control overfitting on imbalanced datasets. Unlike decision trees, they do not encode split-based inductive biases that often suit heterogeneous tabular variables \cite{borisov2022deep}.

Sequence models — LSTMs and GRUs — are natural when transactions form meaningful entity-level sequences with state-dependent fraud patterns. ULB provides an elapsed timestamp but no cardholder identifiers from which such sequences can be reconstructed, whereas BAF represents individual account-opening applications rather than transaction streams. Recurrent models are therefore not a natural match for the structures available in these two benchmarks.

Autoencoder-based anomaly detection trains a reconstruction network on legitimate transactions, deriving anomaly scores from reconstruction error. Under labelled supervised conditions, reconstruction-based approaches often trail discriminative models that optimise the classification objective directly \cite{hilal2022financial}. A recurring challenge for deep learning approaches in fraud detection is sensitivity to extreme class imbalance. Both GBDT and neural implementations can incorporate loss weighting, while data-level resampling changes the training distribution for either family. The interaction between model family and imbalance strategy is therefore treated as an empirical variable in this thesis rather than assumed in advance.

% ============================================================
\section{Transformers for Tabular Fraud Detection}
\label{subsec:transformers_tabular}
% ============================================================

\subsection{The Transformer Architecture}
\label{subsec:sota-transformer-arch}

The Transformer architecture \cite{vaswani2017attention} demonstrated that sequence modelling could be performed without recurrence, enabling parallel processing through attention mechanisms. Its core operation is \textit{multi-head self-attention}: each position attends to other positions through weights learned end-to-end from scaled dot products. The Transformer's success in natural language processing and computer vision motivated extensions to tabular data, a different modality with no inherent sequential or spatial ordering among features.

\subsection{Tabular Transformer Architectures}
\label{subsec:sota-tabular-transformers}

\textbf{TabTransformer} \cite{huang2020tabtransformer} was among the first applications of self-attention to tabular data. Its design applies column-wise attention to categorical feature embeddings, while numerical features bypass the attention blocks and are concatenated with the contextual embeddings before prediction. Consequently, numerical features — all predictors in ULB and most predictors in BAF — receive no attention-based contextual transformation.

\textbf{SAINT} (Self-Attention and Intersample Attention Transformer) \cite{somepalli2021saint} extends attention along two axes: column-wise attention across features within a row, and row-wise attention across samples within a batch. The dual attention mechanism is designed to capture inter-sample dependencies, which may be relevant in fraud detection when groups of transactions share coordinated patterns. SAINT reports competitive performance against GBDT models on several tabular benchmarks, though results are inconsistent across datasets and tasks.

\textbf{FT-Transformer} (Feature Tokenizer + Transformer) \cite{gorishniy2021revisiting} tokenises all features — both categorical and numerical — as dense embedding vectors and applies multi-head self-attention over the full set of feature tokens. A class token is appended for classification, following the vision transformer (\textit{ViT}) convention. A subsequent contribution \cite{gorishniy2022embeddings} reports that replacing linear numerical embeddings with periodic (piecewise-linear or Fourier-basis) embeddings can improve performance on datasets where numerical features dominate — a property relevant to financial fraud datasets such as ULB, where all predictors are continuous. The FT-Transformer is the architecture employed in this thesis as the representative tabular transformer. Its design allows cross-feature interactions to be learned without assuming a natural sequence among columns. Whether that flexibility provides an advantage over the piecewise-constant inductive bias of GBDTs on fraud data is an empirical question rather than an architectural certainty.

\subsection{The GBDT vs.\ Deep Learning Debate on Tabular Data}
\label{subsec:sota-dl-gbdt-debate}

The empirical question of whether tabular transformers outperform GBDTs has generated a structured research debate. Shwartz-Ziv and Armon \cite{shwartzziv2022tabular} conducted a controlled benchmark across multiple tabular datasets and found that GBDT methods — when properly tuned — match or outperform tabular deep learning methods in the majority of cases. They attribute this pattern partly to inductive biases that are well suited to the statistical properties of many tabular datasets, in contrast to self-attention's origins in correlated sequential tokens. Borisov et al. \cite{borisov2022deep} reach a similar conclusion in a comprehensive survey: tabular transformers have closed the gap with GBDTs on some benchmarks, but GBDTs retain an advantage on many heterogeneous, small-to-medium tabular datasets, and no single deep learning architecture consistently outperforms well-tuned GBDTs.

In the specific domain of financial fraud, Thimonier et al. \cite{thimonier2024comparative} report higher LightGBM performance than the evaluated anomaly-detection baselines on credit card fraud data. On the BAF dataset, Sun et al. \cite{sun2025objective} report that objective-function choice has a larger effect than the evaluated architectural differences, providing further evidence against assuming systematic transformer gains over optimised alternatives.

Several hypotheses may explain why GBDTs often retain an advantage on financial tabular data. Their split-based inductive bias can represent irregular decision boundaries efficiently, whereas attention-based models must estimate a larger set of continuous parameters and may require more data and regularisation. Probability calibration can also differ across model families and datasets, although it is not guaranteed by either architecture. Finally, synthetic oversampling may interact differently with tree partitioning and gradient-based optimisation, particularly when numerical and categorical variables are represented jointly. These are plausible mechanisms to be evaluated cautiously; aggregate performance alone cannot identify which mechanism is responsible.

Whether these theoretical considerations translate into consistent empirical gaps on fraud-specific benchmarks is an open question that this thesis addresses empirically. The accompanying diagnostics are used to characterise fitted-model behaviour, not to identify a causal mechanism from aggregate performance alone.

% ============================================================
\section{Explainable Artificial Intelligence in Fraud Detection}
\label{subsec:XAI}
% ============================================================

\subsection{Motivation: Compliance, Trust, and Operational Constraints}
\label{subsec:sota-xai-motivation}

The deployment of machine learning systems in financial decision-making is increasingly subject to transparency, governance, and accountability expectations. Regulations such as the European General Data Protection Regulation (GDPR), together with national banking supervisory frameworks, motivate documented oversight of automated decisions, although the precise legal obligations depend on the decision context and system use \cite{nesvijevskaia2021accuracy}. Beyond regulatory compliance, operational teams benefit from explanations when auditing model decisions, identifying systematic errors, and maintaining justified confidence in automated systems.

The trade-off between predictive accuracy and interpretability is well-documented \cite{molnar2022interpretable}. Logistic regression and shallow decision trees permit relatively direct inspection of coefficients or splits, but may be less expressive than ensemble and deep models. Post-hoc explanation methods address this tension by providing local or global attributions of black-box model decisions without requiring inspection of model internals.

\subsection{SHAP: Shapley Additive Explanations}
\label{subsec:sota-shap}

SHAP \cite{lundberg2017unified} is a widely adopted post-hoc explanation framework in fraud detection research. It grounds feature attributions in cooperative game theory: the SHAP value of feature $j$ for instance $i$ is based on the average marginal contribution of feature $j$ across feature coalitions. Within the additive explanation framework and the assumptions of the selected explainer, the resulting attributions satisfy properties including local accuracy, missingness, and consistency.

For tree-based models, TreeSHAP \cite{lundberg2020local} computes exact values for its tree-specific SHAP formulation in polynomial time, making it computationally feasible for large ensembles. For neural networks, model-agnostic or gradient-based approximations are commonly used, at greater computational cost and with additional approximation assumptions. This asymmetry is practically significant: the tree-specific computation is efficient, whereas neural-model explanations are typically approximate and implementation-dependent.

A common conflation in the deep learning literature is treating attention weights as feature importance. Attention weights indicate which tokens a given position attends to during the forward pass, but they do not constitute a faithful attribution of model output to input features in the Shapley value sense \cite{molnar2022interpretable}. SHAP must be applied post-hoc as a separate computation even when attention is available.

\subsection{Interpretability Consistency Across Model Families}
\label{subsec:sota-xai-cross-model}

Whether different model families identify the same features as important — cross-model interpretability consistency — is understudied in the fraud detection literature. Walauskis and Khoshgoftaar \cite{walauskis2024scalable} apply SHAP-based feature selection to large-scale tabular data but do not compare SHAP ranking consistency across model types. Models with similar aggregate PR-AUC may still generate different feature attributions; such disagreement is relevant evidence for model governance, although attribution agreement alone cannot establish that either model has learned the correct reasoning process.

LIME \cite{ribeiro2016lime} provides an alternative post-hoc explanation approach: for a given instance, it fits a locally linear approximation to the model's predictions in a neighbourhood of that instance. LIME attributions are approximate by construction and can vary with the sampled neighbourhood. SHAP is commonly selected when additive attributions and a consistent comparison framework are desired, but its interpretation still depends on the explainer, background distribution, and feature-dependence assumptions \cite{molnar2022interpretable}.

% ============================================================
\section{Synthesis, Limitations of Existing Work, and Research Gaps}
\label{subsec:limitations_research}
% ============================================================

\subsection{Recurring Methodological Limitations}
\label{subsec:sota-limitations}

The literature reviewed in this chapter reveals four recurring methodological problems that can compromise the reliability and comparability of published results.

\textbf{Data leakage via incorrect resampling.} A consequential error is the application of oversampling before the train/test split, or within cross-validation without isolating resampling to the training portion of each fold. Hayat et al. \cite{hayat2025leakage} demonstrate empirically that this practice can inflate reported PR-AUC and $F_1$ scores sufficiently to alter the ranking of competing methods. Despite being well understood in the methodology literature, the error persists in fraud detection studies.

\textbf{Inappropriate evaluation metrics.} Using ROC-AUC as the sole primary metric, as detailed in Section~\ref{subsec:fraud_detection_class_imbalance}, can be inadequate under extreme class imbalance \cite{saito2015precision, imani2026why}. Studies that report only ROC-AUC or accuracy provide an incomplete characterisation of fraud detection capability.

\textbf{Ad hoc threshold selection.} The decision threshold is rarely treated as a principled experimental variable. Papers that fix $\tau = 0.5$ without justification, or that optimise $\tau$ on the test set, introduce evaluation bias that conflates threshold selection with model quality. Probability calibration, which is required when scores are interpreted as posterior risks or linked to explicit cost ratios, is equally rarely discussed.

\textbf{Single-dataset studies without robustness assessment.} Many published comparisons use the ULB dataset alone. While ULB is a valuable benchmark, its two-day temporal window, PCA-transformed features, and specific imbalance ratio limit the generalisability of conclusions drawn from it. The BAF suite, designed explicitly for distribution-shift evaluation, remains less commonly used in comparative model studies.

\subsection{Research Gaps and Positioning of This Work}
\label{subsec:sota-gaps}

The preceding limitations identify specific gaps that this thesis addresses through six experimental contributions.

\textit{Gap 1: Limited joint, leakage-free comparison of GBDTs and tabular transformers on fraud data.} To the best of our knowledge, no published study provides a simultaneous, controlled comparison of classical gradient-boosted models and the FT-Transformer on both ULB and BAF under a shared leakage-free evaluation framework. The tabular benchmarks of Shwartz-Ziv and Armon \cite{shwartzziv2022tabular} and Borisov et al. \cite{borisov2022deep} suggest GBDTs remain competitive, but fraud datasets differ structurally from general tabular benchmarks through extreme imbalance, regulated use, and behavioural and device features.

\textit{Gap 2: Threshold selection is rarely treated as a first-class experimental variable.} To the best of our knowledge, existing fraud detection studies do not systematically evaluate fixed, max-F1, max-F2, and precision-constrained rules across the same model families and both benchmark datasets under a leakage-free protocol.

\textit{Gap 3: Factorial interactions between imbalance strategy and model type are understudied.} The existing literature commonly treats imbalance strategies and model choice as separate decisions. To the best of our knowledge, a full factorial design crossing the seven selected strategies with five supervised models on ULB and BAF has not been reported.

\textit{Gap 4: Cross-domain generalisation on BAF variants is insufficiently characterised across model families.} Understanding how GBDTs, Random Forest, linear, one-class, and transformer models respond to the controlled group-size, prevalence, separability, and temporal conditions in the BAF suite has direct operational relevance but remains insufficiently explored.

\textit{Gap 5: SHAP consistency across compatible model explainers and dataset variants is insufficiently documented.} Whether linear and tree-based models that achieve similar PR-AUC identify overlapping top-feature sets, and whether the Base-trained LGBM set remains stable across BAF Variant test distributions, is insufficiently documented.

\textit{Gap 6: FT-Transformer attention patterns on fraud data are insufficiently characterised.} In particular, the extent to which the model's last-layer attention is concentrated or diffuse across BAF feature tokens has received limited analysis. Attention is treated here as a descriptive architectural diagnostic, not as a faithful substitute for feature attribution.

These six gaps map directly to the experimental contributions detailed in Chapter~\ref{sec:chapterMethodology} (p.~\pageref{sec:chapterMethodology}). Before presenting their implementation, Chapter~\ref{sec:chapterTheoreticalBackground} (p.~\pageref{sec:chapterTheoreticalBackground}) defines the algorithms, imbalance interventions, metrics, threshold rules, and explanation methods required to interpret the experimental design.
